\chapter{2008年浙江卷（理）}

\section{单选题}

\begin{problem}
已知 \(a\) 是实数，\( \frac{a-\i}{1+\i} \) 是纯虚数，则 \(a=\)（\quad）

\choices
    {1}
    {\(-1\)}
    {\(\sqrt{2}\)}
    {\(-\sqrt{2}\)}
\end{problem}

\begin{answer}
A
\end{answer}

\begin{solution}
将分式的分子、分母同乘以共轭复数 \(1-\i\)，得
\[
\frac{a-\i}{1+\i}=\frac{(a-\i)(1-\i)}{2}=\frac{a-1}{2}-\frac{a+1}{2}\i
\]
它是纯虚数，所以实部为零，即 \(a-1=0\)，从而 \(a=1\)，选 A.
\end{solution}

\begin{problem}
已知 \(U = \R\)，\(A = \{x \mid x > 0\}\)，\(B = \{x \mid x \leqslant -1\}\)，则 \(\left(A \cap \complement_U B\right) \cup \left(B \cap \complement_U A\right) =\)（\quad）

\choices
    {\(\varnothing\)}
    {\(\{x\mid x\leqslant 0\}\)}
    {\( \{x \mid x > -1\} \)}
    {\(\{x\mid x > 0\text{或} x\leqslant -1\}\)}
\end{problem}

\begin{answer}
D
\end{answer}

\begin{solution}
因为 \(B=\{x\mid x\leqslant-1\}\)，所以 \(\complement_U B=\{x\mid x>-1\}\).于是
\(A\cap\complement_U B=A=\{x\mid x>0\}\).又因为 \(\complement_U A=\{x\mid x\leqslant0\}\)，所以
\(B\cap\complement_U A=B=\{x\mid x\leqslant-1\}\).两部分的并集为
\(\{x\mid x>0\text{或}x\leqslant-1\}\)，选 D.
\end{solution}

\begin{problem}
已知 \(a\)，\(b\) 都是实数，那么“\(a^2 > b^2\)”是“\(a > b\)”的（\quad）

\choices
    {充分而不必要条件}
    {必要而不充分条件}
    {充分必要条件}
    {既不充分也不必要条件}
\end{problem}

\begin{answer}
D
\end{answer}

\begin{solution}
若 \(a^2>b^2\)，不一定有 \(a>b\)，例如 \(a=-2\)，\(b=1\) 时，\(a^2>b^2\) 成立而 \(a>b\) 不成立；故前者不是后者的充分条件.
若 \(a>b\)，也不一定有 \(a^2>b^2\)，例如 \(a=1\)，\(b=-2\) 时，\(a>b\) 成立而 \(a^2>b^2\) 不成立；故前者不是后者的必要条件.因此它是既不充分也不必要条件，选 D.
\end{solution}

\begin{problem}
在 \((x - 1)(x - 2)(x - 3)(x - 4)(x - 5)\) 的展开式中，含 \(x^4\) 的项的系数是（\quad）

\choices
    {\(-15\)}
    {85}
    {\(-120\)}
    {274}
\end{problem}

\begin{answer}
A
\end{answer}

\begin{solution}
要取得 \(x^4\) 项，必须从五个因式中选取一个常数项，其余四个因式均选取 \(x\).因此该项的系数为
\(-1-2-3-4-5=-15\)，选 A.
\end{solution}

\begin{problem}
在同一平面直角坐标系中，函数 \( y = \cos \left( \frac{x}{2} + \frac{3\pi}{2} \right) \)（\(x \in [0, 2\pi]\)）的图象和直线 \( y = \frac{1}{2} \) 的交点个数是（\quad）

\choices
    {0}
    {1}
    {2}
    {4}
\end{problem}

\begin{answer}
C
\end{answer}

\begin{solution}
令 \(u=\frac{x}{2}\)，则 \(u\in[0,\pi]\)，且
\(\cos\left(\frac{x}{2}+\frac{3\pi}{2}\right)=\sin\frac{x}{2}=\sin u\).
方程 \(\sin u=\frac12\) 在 \([0,\pi]\) 内的解为 \(u=\frac\pi6\)，\(\frac{5\pi}6\)，对应
\(x=\frac\pi3\)，\(\frac{5\pi}3\).所以交点有 2 个，选 C.
\end{solution}

\begin{problem}
已知 \(\{a_{n}\}\) 是等比数列，\(a_{2} = 2\)，\(a_{5} = \frac{1}{4}\)，则 \(a_{1}a_{2} + a_{2}a_{3} + \cdots + a_{n}a_{n+1} =\)（\quad）

\choices
    {\(16(1 - 4^{-n})\)}
    {\(16(1 - 2^{-n})\)}
    {\(\frac{32}{3}(1 - 4^{-n})\)}
    {\(\frac{32}{3}(1 - 2^{-n})\)}
\end{problem}

\begin{answer}
C
\end{answer}

\begin{solution}
设等比数列的公比为 \(q\).由已知条件，
\(a_5=a_2q^3\)，所以 \(2q^3=\frac14\)，即 \(q=\frac12\).于是
\(a_n=2\left(\frac12\right)^{n-2}=2^{3-n}\)，从而
\(a_na_{n+1}=2^{5-2n}=8\left(\frac14\right)^{n-1}\).因此
\[
\sum_{k=1}^{n}a_ka_{k+1}=8\frac{1-\left(\frac14\right)^n}{1-\frac14}=\frac{32}{3}(1-4^{-n})
\]
所以选 C.
\end{solution}

\begin{problem}
若双曲线 \(\frac{x^2}{a^2} - \frac{y^2}{b^2} = 1\) 的两个焦点到一条准线的距离之比为 3:2，则双曲线的离心率是（\quad）

\choices
    {3}
    {5}
    {\( \sqrt{3} \)}
    {\( \sqrt{5} \)}
\end{problem}

\begin{answer}
D
\end{answer}

\begin{solution}
设双曲线的焦半距为 \(c\)，离心率为 \(e=\frac ca>1\).取右准线
\(x=\frac ae\)，则两个焦点到该准线的距离分别为
\(c-\frac ae\) 和 \(c+\frac ae\).较大者与较小者之比为 \(\frac32\)，故
\[
\frac{c+\frac ae}{c-\frac ae}=\frac32
\]
代入 \(c=ae\)，得 \(\frac{e^2+1}{e^2-1}=\frac32\)，从而 \(e^2=5\).由于 \(e>0\)，所以 \(e=\sqrt5\)，选 D.
\end{solution}

\begin{problem}
若 \(\cos \alpha + 2\sin \alpha = -\sqrt{5}\)，则 \(\tan \alpha =\)（\quad）

\choices
    {\( \frac{1}{2} \)}
    {2}
    {\(-\frac{1}{2}\)}
    {\(-2\)}
\end{problem}

\begin{answer}
B
\end{answer}

\begin{solution}
由柯西不等式，
\(\cos\alpha+2\sin\alpha\geqslant-\sqrt{1^2+2^2}\sqrt{\cos^2\alpha+\sin^2\alpha}=-\sqrt5\).
题设取到等号，所以单位向量 \((\cos\alpha,\sin\alpha)\) 必与 \((1,2)\) 反向，即
\(\cos\alpha=-\frac1{\sqrt5}\)，\(\sin\alpha=-\frac2{\sqrt5}\).因此
\(\tan\alpha=\frac{\sin\alpha}{\cos\alpha}=2\)，选 B.
\end{solution}

\begin{problem}
已知 \(\bs{a}\)，\(\bs{b}\) 是平面内两个互相垂直的单位向量，若向量 \(\bs{c}\) 满足 \((\bs{a} - \bs{c}) \cdot (\bs{b} - \bs{c}) = 0\)，则 \(|\bs{c}|\) 的最大值是（\quad）

\choices
    {1}
    {2}
    {\(\sqrt{2}\)}
    {\(\frac{\sqrt{2}}{2}\)}
\end{problem}

\begin{answer}
C
\end{answer}

\begin{solution}
以 \(\bs{a}\)，\(\bs{b}\) 为正交单位基底，设 \(\bs{c}=x\bs{a}+y\bs{b}\).由题设
\[
(\bs{a}-\bs{c})\cdot(\bs{b}-\bs{c})=\bs{c}\cdot\bs{c}-\bs{a}\cdot\bs{c}-\bs{b}\cdot\bs{c}=x^2+y^2-x-y=0
\]
所以 \(\left(x-\frac12\right)^2+\left(y-\frac12\right)^2=\frac12\)，即点 \((x,y)\) 在圆心
\(\left(\frac12,\frac12\right)\)、半径 \(\frac{\sqrt2}{2}\) 的圆上.原点到圆心的距离也是 \(\frac{\sqrt2}{2}\)，故
\(|\bs{c}|\) 的最大值为两者之和 \(\sqrt2\)，选 C.
\end{solution}

\begin{problem}
如图，\(AB\) 是平面 \(\alpha\) 的斜线段，\(A\) 为斜足. 若点 \(P\) 在平面 \(\alpha\) 内运动，使得 \(\triangle ABP\) 的面积为定值，则动点 \(P\) 的轨迹是（\quad）

\bitmapfigure[max width=0.44\linewidth]{img/2008/zhejiang_science/q10_fig1.png}

\choices
    {圆}
    {椭圆}
    {一条直线}
    {两条平行直线}
\end{problem}

\begin{answer}
B
\end{answer}

\begin{solution}
设点 \(B\) 在平面 \(\alpha\) 上的正投影为 \(H\)，则 \(BH\perp\alpha\)，且 \(BH>0\)，因为 \(AB\) 是斜线段.
设平面 \(\alpha\) 内以 \(A\) 为原点、\(AH\) 为 \(x\) 轴，垂直于 \(AH\) 的方向为 \(y\) 轴，并记 \(AH=d>0\)、\(BH=h>0\). 对动点 \(P=(x,y)\)，有
\[
4S_{\triangle ABP}^{\,2}=|\overrightarrow{AB}\times\overrightarrow{AP}|^2
 =h^2(x^2+y^2)+d^2y^2=h^2x^2+(h^2+d^2)y^2
\]
由于面积为定值且为正，设 \(4S_{\triangle ABP}^{\,2}=4S^2>0\)，则
\[
\frac{x^2}{4S^2/h^2}+\frac{y^2}{4S^2/(h^2+d^2)}=1
\]
这是以 \(A\) 为中心的椭圆，所以选 B.
\end{solution}

\section{填空题}

\begin{problem}
已知 \(a > 0\)，若平面内三点 \(A(1, -a)\)，\(B(2, a^2)\)，\(C(3, a^3)\) 共线，则 \(a =\)\fillinblank{}.
\end{problem}

\begin{answer}
\(1+\sqrt{2}\)
\end{answer}

\begin{solution}
三点共线，所以斜率相等.由 \(AB\) 与 \(AC\) 的斜率分别为
\(a^2+a\) 和 \(\frac{a^3+a}{2}\)，得
\[
a^2+a=\frac{a^3+a}{2}
\]
即 \(a(a^2-2a-1)=0\).因为 \(a>0\)，故
\(a=1+\sqrt2\) 或 \(a=1-\sqrt2<0\)，舍去后者，得到 \(a=1+\sqrt2\).
\end{solution}

\begin{problem}
已知 \(F_1\)，\(F_2\) 为椭圆 \(\frac{x^2}{25} + \frac{y^2}{9} = 1\) 的两个焦点，过 \(F_1\) 的直线交椭圆于 \(A\)，\(B\) 两点. 若 \(|F_2A| + |F_2B| = 12\)，则 \(|AB| =\fillinblank{}\).
\end{problem}

\begin{answer}
\(8\)
\end{answer}

\begin{solution}
椭圆的长半轴为 \(5\)，短半轴为 \(3\)，所以焦半距
\(c=\sqrt{25-9}=4\).对椭圆上任一点 \(X\)，有 \(|F_1X|+|F_2X|=10\).
直线过椭圆内部的焦点 \(F_1\)，故 \(F_1\) 在线段 \(AB\) 上，从而
\(|AB|=|F_1A|+|F_1B|\).于是
\[
12=|F_2A|+|F_2B|=(10-|F_1A|)+(10-|F_1B|)=20-|AB|
\]
所以 \(|AB|=8\).
\end{solution}

\begin{problem}
在 \(\triangle ABC\) 中，角 \(A\)、\(B\)、\(C\) 所对的边分别为 \(a\)、\(b\)、\(c\)，若 \((\sqrt{3} b - c) \cos A = a \cos C\)，则 \(\cos A =\)\fillinblank{}.
\end{problem}

\begin{answer}
\(\frac{\sqrt{3}}{3}\)
\end{answer}

\begin{solution}
由正弦定理，设公共比例因子为 \(k\)，则
\(a=k\sin A\)，\(b=k\sin B\)，\(c=k\sin C\).原条件除以 \(k\) 后为
\[
(\sqrt3\sin B-\sin C)\cos A=\sin A\cos C
\]
又 \(B=\pi-A-C\)，所以 \(\sin B=\sin(A+C)=\sin A\cos C+\cos A\sin C\).代入并移项，得
\[
(\sqrt3\cos A-1)(\sin A\cos C+\cos A\sin C)=0
\]
即 \((\sqrt3\cos A-1)\sin B=0\).三角形内角满足 \(0<B<\pi\)，故 \(\sin B>0\)，从而
\(\cos A=\frac1{\sqrt3}=\frac{\sqrt3}{3}\).
\end{solution}

\begin{problem}
如图，已知球 \(O\) 的表面上四点 \(A\)、\(B\)、\(C\)、\(D\)，\(DA \perp\) 平面 \(ABC\)，\(AB \perp BC\)，\(DA = AB = BC = \sqrt{3}\)，则球 \(O\) 的体积等于\fillinblank{}.

\bitmapfigure[max width=6.0cm]{img/2008/zhejiang_science/q14_fig1.png}
\end{problem}

\begin{answer}
\(\frac{9\pi}{2}\)
\end{answer}

\begin{solution}
建立直角坐标系，令 \(B=(0,0,0)\)，\(A=(\sqrt3,0,0)\)，\(C=(0,\sqrt3,0)\).由
\(DA\perp\) 平面 \(ABC\) 且 \(DA=\sqrt3\)，可取 \(D=(\sqrt3,0,\sqrt3)\).设球心为
\(O=(u,v,w)\).
因为 \(OA=OB\)，得 \(u=\frac{\sqrt3}{2}\)；因为 \(OC=OB\)，得 \(v=\frac{\sqrt3}{2}\).再由 \(OD=OB\)，
\(2\sqrt3(u+w)=6\)，所以 \(w=\frac{\sqrt3}{2}\).因此
\[
R^2=OB^2=u^2+v^2+w^2=\frac94
\]
得到 \(R=\frac32\).球的体积为
\(\frac43\pi R^3=\frac43\pi\left(\frac32\right)^3=\frac{9\pi}{2}\).
\end{solution}

\begin{problem}
已知 \(t\) 为常数，函数 \(y = |x^2 - 2x - t|\) 在区间 \([0,3]\) 上的最大值为 2 ，则 \(t =\fillinblank{}\).
\end{problem}

\begin{answer}
\(1\)
\end{answer}

\begin{solution}
令 \(h(x)=x^2-2x-t\).在 \([0,3]\) 上，\(h\) 的最小值为
\(h(1)=-1-t\)，最大值为端点中的较大者 \(h(3)=3-t\).因此
\[
\max_{x\in[0,3]}|h(x)|=\max\{|1+t|,|3-t|\}=2
\]
要使该最大值不超过 2，必须同时有 \(-3\leqslant t\leqslant1\) 和 \(1\leqslant t\leqslant5\)，故只能 \(t=1\)，且此时最大值确为 2.
\end{solution}

\begin{problem}
用 \(1\)，\(2\)，\(3\)，\(4\)，\(5\)，\(6\) 组成六位数（没有重复数字），要求任何相邻两个数字的奇偶性不同，且 \(1\) 和 \(2\) 相邻，这样的六位数的个数是\fillinblank{}. （用数字作答）
\end{problem}

\begin{answer}
\(40\)
\end{answer}

\begin{solution}
奇数有 \(1\)，\(3\)，\(5\)，偶数有 \(2\)，\(4\)，\(6\).相邻数字奇偶性不同，所以奇偶位置模式只有两种：
\(OEOEOE\) 和 \(EOEOEO\).

在模式 \(OEOEOE\) 中，数字 \(1\) 位于第 1、3、5 个位置.分别考虑它与数字 \(2\) 相邻时，数字 \(2\) 可占的位置数为 \(1\)，\(2\)，\(2\)，共 \(5\) 种位置安排.其余两个奇数和两个偶数分别有 \(2!\) 和 \(2!\) 种排法，所以共有 \(5\times2!\times2!=20\) 个.另一种模式同样有 \(20\) 个，故总数为 \(40\).
\end{solution}

\begin{problem}
若 \(a \geqslant 0\)，\(b \geqslant 0\)，且当 \(\begin{cases}
x \geqslant 0,\\
y \geqslant 0,\\
x + y \leqslant 1
\end{cases}\) 时，恒有 \(ax + by \leqslant 1\)，则以 \(a\)，\(b\) 为坐标的点 \(P(a, b)\) 所形成的平面区域的面积等于\fillinblank{}.
\end{problem}

\begin{answer}
\(1\)
\end{answer}

\begin{solution}
在三角形区域 \(x\geqslant0\)，\(y\geqslant0\)，\(x+y\leqslant1\) 上，线性函数 \(ax+by\) 的最大值在顶点
\((0,0)\)、\((1,0)\)、\((0,1)\) 中取得.题设等价于
\(a\leqslant1\) 且 \(b\leqslant1\).结合 \(a\geqslant0\)、\(b\geqslant0\)，点 \(P(a,b)\) 的区域为正方形
\(0\leqslant a\leqslant1\)，\(0\leqslant b\leqslant1\)，面积为 \(1\).
\end{solution}

\section{解答题}

\begin{problem}
如图，矩形 \(ABCD\) 和梯形 \(BEFC\) 所在平面互相垂直，\(BE \parallel CF\)，\(\angle BCF = \angle CEF = 90^\circ\)，\(AD = \sqrt{3}\)，\(EF = 2\).

\begin{enumerate}
\item 求证：\(AE \parallel\) 平面 \(DCF\)；
\item 当 \(AB\) 的长为何值时，二面角 \(A - EF - C\) 的大小为 \(60^\circ\)？
\end{enumerate}

\bitmapfigure[max width=0.48\linewidth]{img/2008/zhejiang_science/q18_fig1.png}

\begin{solution}
\begin{enumerate}
\item 设 \(AB=p>0\).建立直角坐标系，令
\(B=(0,0,0)\)，\(C=(0,\sqrt3,0)\)，\(A=(p,0,0)\)，\(D=(p,\sqrt3,0)\)，并令梯形所在平面为 \(x=0\).
由 \(\angle BCF=90^\circ\) 及 \(BE\parallel CF\)，可设
\(E=(0,0,u)\)，\(F=(0,\sqrt3,v)\).令 \(d=v-u\)，由 \(EF=2\) 得
\(3+d^2=4\)，即 \(d^2=1\).又由 \(\angle CEF=90^\circ\)，
\[
\overrightarrow{EC}\cdot\overrightarrow{EF}=(0,\sqrt3,-u)\cdot(0,\sqrt3,d)=3-ud=0
\]
故 \(ud=3\).适当取 \(z\) 轴方向后可令 \(d=1\)，于是 \(u=3\)，即
\(E=(0,0,3)\)，\(F=(0,\sqrt3,4)\).

此时 \(D\)、\(C\)、\(F\) 均满足 \(y=\sqrt3\)，所以平面 \(DCF\) 为 \(y=\sqrt3\).而
\(\overrightarrow{AE}=(-p,0,3)\) 的 \(y\) 分量为零，故 \(AE\parallel\) 平面 \(DCF\).
\item 取平面 \(AEF\) 和平面 \(CEF\) 的法向量
\(\bs{n}_1=\overrightarrow{EA}\times\overrightarrow{EF}=(3\sqrt3,-p,p\sqrt3)\)，\(\bs{n}_2=\overrightarrow{EC}\times\overrightarrow{EF}=(4\sqrt3,0,0)\).
二面角 \(A-EF-C\) 的平面角取锐角，设其为 \(\theta\)，则
\[
\cos\theta=\frac{|\bs{n}_1\cdot\bs{n}_2|}{|\bs{n}_1||\bs{n}_2|}
 =\frac{3\sqrt3}{\sqrt{27+4p^2}}
\]
当 \(\theta=60^\circ\) 时，
\(\frac{3\sqrt3}{\sqrt{27+4p^2}}=\frac12\)，所以
\(27+4p^2=108\)，从而 \(p^2=\frac{81}{4}\).由于 \(p>0\)，得 \(AB=p=\frac92\).
\end{enumerate}
\end{solution}
\end{problem}

\begin{problem}
一个袋中有若干个大小相同的黑球、白球和红球. 已知从袋中任意摸出 1 个球，得到黑球的概率是 \(\frac{2}{5}\)；从袋中任意摸出 2 个球，至少得到 1 个白球的概率是 \(\frac{7}{9}\).

\begin{enumerate}
\item 若袋中共有 10 个球，

\begin{enumerate}
\item 求白球的个数；
\item 从袋中任意摸出 3 个球，记得到白球的个数为 \( \xi \)，求随机变量 \( \xi \) 的数学期望 \( E(\xi) \).
\end{enumerate}

\item 求证：从袋中任意摸出 2 个球，至少得到 1 个黑球的概率不大于 \(\frac{7}{10}\)，并指出袋中哪种颜色的球个数最少.
\end{enumerate}

\begin{solution}
\begin{enumerate}
\item
\begin{enumerate}
\item 设白球有 \(w\) 个.两球中没有白球的概率为 \(1-\frac79=\frac29\)，故
\[
\frac{\mathrm C_{10-w}^{2}}{\mathrm C_{10}^{2}}=\frac29
\]
于是 \((10-w)(9-w)=20\)，解得 \(w=5\) 或 \(w=14\).由 \(0\leqslant w\leqslant10\)，得白球有 \(5\) 个.
\item 设第 \(i\) 次摸到白球的指示变量为 \(I_i\)，则 \(\xi=I_1+I_2+I_3\).每次摸球得到白球的概率为 \(\frac5{10}\)，所以由数学期望的线性性质，
\[
E(\xi)=E(I_1)+E(I_2)+E(I_3)=3\times\frac5{10}=\frac32
\]
\end{enumerate}
\item 设袋中共有 \(N\) 个球，黑、白、红球数分别为 \(K\)，\(W\)，\(R\).由 \(\frac KN=\frac25\)，可设
\(N=5m\)，\(K=2m\)，其中 \(m\) 为正整数.两球中没有白球的概率为 \(\frac29\).令非白球数 \(X=K+R=N-W\)，则
\[
\frac{\mathrm C_{X}^{2}}{\mathrm C_{5m}^{2}}=\frac29,
\qquad\text{即}\qquad X(X-1)=\frac{10m(5m-1)}9
\]
两球中至少有一个黑球的概率为
\[
1-\frac{\mathrm C_{3m}^{2}}{\mathrm C_{5m}^{2}}
\]
又因为 \(m\geqslant1\)，
\[
\frac{\mathrm C_{3m}^{2}}{\mathrm C_{5m}^{2}}\geqslant\frac3{10}
\quad\Longleftrightarrow\quad
30m(3m-1)\geqslant15m(5m-1)
\]
而右边移到左边后为 \(15m(m-1)\geqslant0\).因此至少得到一个黑球的概率不大于 \(\frac7{10}\).

还需比较三种球的个数.若 \(X\geqslant\frac{7m}{2}\)，则
\[
X(X-1)\geqslant\frac{7m}{2}\left(\frac{7m}{2}-1\right)>\frac{10m(5m-1)}9
\]
因为两边之差为 \(\frac{m(241m-86)}{36}>0\)，这与上面的等式矛盾.因此
\(X<\frac{7m}{2}\).于是 \(R=X-2m<\frac{3m}{2}<2m=K\)，并且
\(W-R=7m-2X>0\).所以红球数同时少于黑球数和白球数，红球最少.
\end{enumerate}
\end{solution}
\end{problem}

\begin{problem}
已知曲线 \(C\) 是到点 \(P\left(-\frac{1}{2}, \frac{3}{8}\right)\) 和到直线 \(y = -\frac{5}{8}\) 距离相等的点的轨迹. \(l\) 是过点 \(Q(-1, 0)\) 的直线，\(M\) 是 \(C\) 上（不在 \(l\) 上）的动点，\(A\)、\(B\) 在 \(l\) 上，\(MA \perp l\)，\(MB \perp x\) 轴（如图）.

\begin{enumerate}
\item 求曲线 \(C\) 的方程；
\item 求出直线 \(l\) 的方程，使得 \(\frac{|QB|^2}{|QA|}\) 为常数.
\end{enumerate}

\bitmapfigure[max width=0.42\linewidth]{img/2008/zhejiang_science/q20_fig1.png}

\begin{solution}
\begin{enumerate}
\item 设 \(N(x,y)\) 为曲线 \(C\) 上的点.由距离相等，
\[
\left(x+\frac12\right)^2+\left(y-\frac38\right)^2=\left|y+\frac58\right|^2
\]
展开并约去 \(y^2\)，得 \(x^2+x=2y\)，所以曲线 \(C\) 的方程为
\(y=\frac12(x^2+x)\).
\item 直线 \(l\) 不可能为过 \(Q\) 的竖直直线，否则由 \(B\in l\) 且 \(MB\perp x\) 轴可知 \(M\) 的横坐标为 \(-1\)，而曲线 \(C\) 在该直线上只有点 \(Q\)，与 \(M\notin l\) 矛盾.因此设
\(l:y=k(x+1)\).

设 \(M=(x,f(x))\)，其中 \(f(x)=\frac12x(x+1)\).因为 \(MB\perp x\) 轴且 \(B\in l\)，有
\(B=(x,k(x+1))\)，从而 \(|QB|^2=(1+k^2)(x+1)^2\).取直线 \(l\) 的方向向量 \(\bs{d}=(1,k)\)，设
\(A=Q+\lambda\bs{d}\).由 \(MA\perp l\)，
\[
\lambda=\frac{(M-Q)\cdot\bs{d}}{1+k^2}
 =\frac{(x+1)+kf(x)}{1+k^2}
 =\frac{(x+1)\left(1+\frac{kx}{2}\right)}{1+k^2}
\]
因此
\[
\frac{|QB|^2}{|QA|}=\frac{(1+k^2)^{3/2}|x+1|}{\left|1+\frac{kx}{2}\right|}
\]
要使这个比值对曲线 \(C\) 上的动点 \(M\) 为常数，两个一次式
\(1+\frac{kx}{2}\) 与 \(x+1\) 必须成比例.比较常数项和 \(x\) 项，得
\(1=\frac{k}{2}\)，即 \(k=2\).此时 \(l:y=2(x+1)\)，即
\(l:2x-y+2=0\)，且比值恒为 \(5\sqrt5\).
\end{enumerate}
\end{solution}
\end{problem}

\begin{problem}
已知 \(a\) 是实数，函数 \(f(x) = \sqrt{x}(x - a)\).

\begin{enumerate}
\item 求函数 \(f(x)\) 的单调区间；
\item 设 \(g(a)\) 为 \(f(x)\) 在区间 \([0, 2]\) 上的最小值.

\begin{enumerate}
\item 写出 \(g(a)\) 的表达式；
\item 求 \(a\) 的取值范围，使得 \(-6 \leqslant g(a) \leqslant -2\).
\end{enumerate}
\end{enumerate}

\begin{solution}
\begin{enumerate}
\item 函数定义域为 \([0,+\infty)\).当 \(x>0\) 时，
\[
f'(x)=\frac{3x-a}{2\sqrt{x}}
\]
当 \(a\leqslant0\) 时，\(f'(x)>0\)，所以 \(f(x)\) 在 \([0,+\infty)\) 上单调递增；当 \(a>0\) 时，\(f'(x)<0\) 对
\(0<x<\frac a3\) 成立，\(f'(x)>0\) 对 \(x>\frac a3\) 成立，所以 \(f(x)\) 在
\(\left[0,\frac a3\right]\) 上单调递减，在 \(\left[\frac a3,+\infty\right)\) 上单调递增.
\item
\begin{enumerate}
\item 由上题的单调性，
\[
g(a)=
\begin{cases}
0,&a\leqslant0,\\
-\dfrac{2a}{3}\sqrt{\dfrac a3},&0<a<6,\\
\sqrt2(2-a),&a\geqslant6.
\end{cases}
\]
\item 当 \(a\leqslant0\) 时，\(g(a)=0\)，不满足条件.当 \(0<a<6\) 时，
\(-6\leqslant-\frac{2a}{3}\sqrt{\frac a3}\leqslant-2\) 给出 \(3\leqslant a<6\).
当 \(a\geqslant6\) 时，
\(-6\leqslant\sqrt2(2-a)\leqslant-2\) 给出
\(6\leqslant a\leqslant2+3\sqrt2\).合并两种情形，得
\(3\leqslant a\leqslant2+3\sqrt2\).
\end{enumerate}
\end{enumerate}
\end{solution}
\end{problem}

\begin{problem}
已知数列 \(\{a_{n}\}\)，\(a_{n} \geqslant 0\)，\(a_{1} = 0\)，\(a_{n+1}^{2} + a_{n+1} - 1 = a_{n}^{2}\)（\(n \in \N^*\)）. 记 \(S_{n} = a_{1} + a_{2} + \cdots + a_{n}\)，\(T_{n} = \frac{1}{1 + a_{1}} + \frac{1}{(1 + a_{1})(1 + a_{2})} + \cdots + \frac{1}{(1 + a_1)(1 + a_2)\cdots(1 + a_n)}\). 求证：当 \(n\in \N^*\) 时，

\begin{enumerate}
\item \(a_{n} < a_{n + 1}\)；
\item \(S_{n} > n - 2\)；
\item \(T_{n} < 3\).
\end{enumerate}

\begin{solution}
\begin{enumerate}
\item 由 \(a_n\geqslant0\)，递推式中的 \(a_{n+1}\) 取非负根，故
\[
a_{n+1}=\frac{-1+\sqrt{5+4a_n^2}}2
\]
先证 \(0\leqslant a_n<1\).\(a_1=0<1\)，若 \(0\leqslant a_n<1\)，则
\(5+4a_n^2<9\)，从而 \(0<a_{n+1}<1\)，归纳可得对一切 \(n\) 都成立.
又因为两边均非负，
\[
a_{n+1}>a_n\Longleftrightarrow\sqrt{5+4a_n^2}>2a_n+1
\Longleftrightarrow 5+4a_n^2>(2a_n+1)^2
\Longleftrightarrow a_n<1
\]
所以 \(a_n<a_{n+1}\).
\item \(a_2\) 是方程 \(x^2+x-1=0\) 的非负根，故
\(a_2=\frac{\sqrt5-1}{2}>\frac12\).当 \(n=1\)，\(2\) 时，分别有
\(S_1=0>-1\) 和 \(S_2=a_2>0\).

当 \(k\geqslant3\) 时，递推式可改写为
\[
1-a_k=a_k^2-a_{k-1}^2=(a_k-a_{k-1})(a_k+a_{k-1})
\]
由于数列递增且 \(a_2>\frac12\)，有 \(a_k+a_{k-1}>2a_2>1\)，所以
\(1-a_k<a_k-a_{k-1}\).于是当 \(n\geqslant3\) 时，
\[
n-S_n=\sum_{k=1}^{n}(1-a_k)
<1+(1-a_2)+(a_n-a_2)<3-2a_2<2
\]
故 \(S_n>n-2\).
\item 由 \(a_1=0\) 且数列递增，\(a_k\geqslant a_2\) 对 \(k\geqslant2\) 成立.因此
\[
\prod_{j=1}^{k}(1+a_j)\geqslant(1+a_2)^{k-1}
\]
从而
\[
T_n\leqslant\sum_{k=1}^{n}\frac1{(1+a_2)^{k-1}}
<\sum_{k=1}^{\infty}\frac1{(1+a_2)^{k-1}}
=\frac{1+a_2}{a_2}=\frac{3+\sqrt5}{2}<3
\]
\end{enumerate}
\end{solution}
\end{problem}
